\documentclass[11pt]{article}
\usepackage{fontspec,amsmath,amssymb,amsthm,hyperref}
\usepackage[margin=1in]{geometry}
\setlength{\emergencystretch}{3em}
\title{The exact averaged-smoothness one-step LMC theorem}
\author{LangevinLean mathematical exposition}
\date{}
\begin{document}\maketitle
\noindent Source module: \texttt{AverageLMC/GaussianLMC.lean}. Compile this proof document with XeLaTeX or LuaLaTeX. The accompanying HTML pairs the mathematical argument with the exact Lean statements and records the scope of the formalization.


This is Theorem 1, claim 1, for the actual Gaussian Markov transition and the literal average coordinate Taylor supremum. It assumes no coupling recurrence, stationary-reference bundle, or desired convergence estimate.



For the actual independent Gaussian step, \(M_{\rm av}\) is the paper's average of the exact coordinate constants, and \(L\) is the global gradient Lipschitz constant. Theorem 1's local inequality is


\[W_2^2(\nu^+,\pi)\le\frac{1-mh}{1+mh}W_2^2(\nu,\pi)+\frac{(M_{\rm av}+m)h^2p}{1+mh}+\frac{h^2(hL-1)}{1+mh}\int\|\nabla f(x)\|^2\,d\nu(x),\qquad h\ge0.\]
\section{Include the zero step exactly}


At h=0, the deterministic displacement and Gaussian coefficient both vanish. The transition is the projection onto the first coordinate of \ensuremath{\mu} times the standard Gaussian probability law; its pushforward is \ensuremath{\mu}. This proves the zero-step case without a limiting argument.



\noindent Lean declarations: \texttt{ulaLawStep\_zero}.

\section{Use exact coordinate smoothing and take the transport infimum}


Instantiate the canonical coupling theorem with the proved standard-Gaussian coordinate smoothing bound and Mav equal to the average of the literal coordinate Taylor suprema. Every finite initial coupling satisfies the same quadratic estimate. The squared-transport infimum theorem passes this estimate to the true initial Wasserstein distance; it also handles the possible negative coefficient 1−mh when h is large. Divide by the strictly positive 1+mh and collect the terms to obtain the paper’s exact formula for every h≥0. The finite-distance hypothesis expresses the finite-second-moment interpretation of this displayed real-valued inequality.



\noindent Lean declarations: \texttt{lmc\_one\_step}.

\end{document}
